\clearpage
\fancyhead[L]{\small COMET.}
\hypertarget{native-442-part-lower}{}
\begin{center}
{\large COMET.}

\noindent\rule{\linewidth}{0.4pt}

Göttingische gelehrte Anzeigen. Stück 78. Seite 769..771. 1821 Mai 17.

\noindent\rule{\linewidth}{0.4pt}
\end{center}

Wir haben mit der Anzeige der sämmtlichen auf hiesiger Sternwarte gemachten Beobachtungen
des diesjährigen Cometen bis jetzt gezögert, weil wir hofften, dass die auf dieselben gegründeten Resultate durch viel spätere auswärtige Beobachtungen noch würden vervollkommnet werden können.
Da diese Hoffnung unerfüllt geblieben ist, so theilen wir nunmehro die sämmtlichen hiesigen Beobachtungen hier mit.

\Needspace{14\baselineskip}
\SourceLeaf{443}
\begin{center}
\begin{tabular}{lr|r|r|r}
\multicolumn{2}{c|}{1821} & Mittlere Zeit & Gerade Aufst. & Nordl. Abw. \\
\hline
Jan. & 30 & $7^{\mathrm{u}}\;34^{\mathrm{m}}\;32^{\mathrm{s}}$ & $359^\circ\;27'\;7''$ & $16^\circ\;4'\;36''$ \\
Febr. & 3 & $7\quad3\quad56$ & $359\quad3\quad54$ & $15\quad46\quad3$ \\
& 7 & $6\quad42\quad38$ & $358\quad45\quad5$ & $15\quad29\quad49$ \\
& 9 & $6\quad42\quad30$ & $358\quad36\quad24$ & $15\quad21\quad22$ \\
& 10 & $6\quad52\quad27$ & $358\quad32\quad19$ & $16\quad17\quad50$ \\
& 11 & $7\quad12\quad0$ & $358\quad28\quad27$ & $15\quad14\quad26$ \\
März & 1 & $7\quad18\quad7$ & $357\quad13\quad34$ & $14\quad8\quad37$ \\
& 5 & $7\quad5\quad28$ & $356\quad54\quad11$ & $13\quad43\quad5$ \\
\end{tabular}
\end{center}

\Needspace{14\baselineskip}
Hr von \textsc{Staudt}, von dessen ausgezeichneter Geschicklichkeit im astronomischen Calcül wir schon öfters Proben mitgetheilt haben, gründet auf diese Beobachtungen folgende \emph{parabolische Elemente}.

\begin{center}
\small
\begin{tabular}{p{.73\linewidth}r}
Durchgangszeit durch die Sonnennähe 1821 März \mbox{.\,.\,.\,.\,.\,.\,.\,.\,.} & $21.61890$ \\
Logarithm des kleinsten Abstandes von der Sonne Meridian von Göttingen & $8.9645990$ \\
Argument der Breite in der Sonnennähe \mbox{.\,.\,.\,.\,.\,.\,.\,.\,.\,.\,.\,.\,.} & $169^\circ\;10'\;9''3$ \\
Länge des aufsteigenden Knoten \mbox{.\,.\,.\,.\,.\,.\,.\,.\,.\,.\,.\,.\,.\,.\,.\,.} & $48\quad44\quad14.7$ \\
Neigung der Bahn \mbox{.\,.\,.\,.\,.\,.\,.\,.\,.\,.\,.\,.\,.\,.\,.\,.\,.\,.\,.\,.\,.} & $106\quad40\quad16.4$ \\
\end{tabular}
\end{center}

\Needspace{17\baselineskip}
Diese Elemente stellen die sämmtlichen uns bekannt gewordenen Beobachtungen so schön dar, dass keine weitere Verbesserung, und noch weniger eine Bestimmung der Ellipticität, gemacht werden kann, wenn nicht noch viel spätere gute Beobachtungen von andern Orten bekannt werden. Wir können des beschränkten Raums wegen hier nur die Vergleichung der Elemente mit den hiesigen Beobachtungen beifügen:

\begin{center}
\begin{tabular}{lr|r|r}
& & \multicolumn{2}{c}{Unterschied.} \\
& & Ger. Aufst. & Abweich. \\
\hline
Januar & 30 & $-18''8$ & $+34''9$ \\
Februar & 3 & $+3.9$ & $-6.0$ \\
& 7 & $-13.1$ & $-36.0$ \\
& 9 & $-0.1$ & $+14.7$ \\
& 10 & $+2.6$ & $+7.3$ \\
& 11 & $-1.2$ & $-3.5$ \\
März & 1 & $-0.3$ & $-0.3$ \\
& 5 & $-5.3$ & $-20.6$ \\
\end{tabular}
\end{center}

\Needspace{11\baselineskip}
Dieser Comet zeichnet sich durch seine grosse Annäherung an die Sonne vor andern aus, und musste daher um die Zeit des Perihelium eine sehr grosse Lichtstärke erreichen. Das Licht des Cometen als erborgtes Sonnenlicht betrachtet, gab Hrn. \textsc{von Staudts} Rechnung die Lichtstärke:

\begin{center}
\begin{tabular}{lr|r}
Januar & 21 & $0.139$ \\
Februar & 10 & $0.258$ \\
März & 1 & $0.832$ \\
& 5 & $1.253$ \\
\end{tabular}
\end{center}
\Needspace{7\baselineskip}
\begin{center}
\begin{tabular}{lr|r}
März & 18 & $19.588$ \\
& 19 & $32.519$ \\
& 20 & $59.932$ \\
& 21 & $110.610$ \\
\end{tabular}
\end{center}
\Needspace{7\baselineskip}
\begin{center}
\begin{tabular}{lr|r}
März & 22 & $128.139$ \\
& 23 & $81.238$ \\
& 24 & $46.005$ \\
& 25 & $27.980$ \\
\end{tabular}
\end{center}

Die Zahlen vom 18. März an gelten für die Culminationszeit in Göttingen. In den ersten Tagen
des März konnte der Comet sehr bequem mit blossen Augen gesehen werden und glich an Helligkeit
fast einem Sterne 3ter Grösse. Es war daher allerdings die grösste Hoffnung, dass der Comet zur Zeit
seiner grössten Helligkeit mit lichtstarken Instrumenten trotz seiner nahen Stellung bei der Sonne, bei
Tage im Meridian würde beobachtet werden können. Auf der hiesigen Sternwarte erlaubte inzwischen
der bedeckte Himmel vom 19. bis 23. März keinen Versuch; am 18. und 24. hingegen, wo um die 
\SourceLeaf{444}%
Culminationszeit des Cometen der Himmel klar war, konnte keine Spur des Cometen bemerkt werden, obgleich dessen Ort genau voraus berechnet war. Nach den Nachrichten, welche Hr. Hofr. \textsc{Gauss} von
verschiedenen andern Sternwarten erhalten hat, sind daselbst die Nachsuchungen um die Zeit der grössten Helligkeit des Cometen gleichfalls durch ungünstige Witterung vereitelt worden.
